\chapter{Elementary Aspects of Metric Geometry}
\label{app:chow-iii-elementary-metric-geometry}
\dcref{appG}
\source{chow-et-al-ricci-flow-part-iii:page-0408}

\newcommand{\GH}{d_{\mathrm{GH}}}
\newcommand{\GHpt}{d^{\mathrm{pt}}_{\mathrm{GH}}}
\newcommand{\diam}{\operatorname{diam}}
\newcommand{\Cone}{\operatorname{Cone}}
\newcommand{\cL}{\mathcal L}
\newcommand{\cA}{\mathcal A}
\newcommand{\R}{\mathbb R}
\newcommand{\dH}{d_{\mathrm H}}
\newcommand{\dimH}{\operatorname{dim}_{\mathrm H}}

This appendix records the metric, length-space, and Aleksandrov-space
background used in the Ricci-flow chapters.

\section{Metric spaces and length spaces}

\begin{remark}[Scaling metric spaces]
\label{rem:chow-iii-g1-scaling-metric-spaces}
\source{chow-et-al-ricci-flow-part-iii:page-0410}
\dcref{appG:G.1}\group{Metric spaces}
If $(X,d)$ is a metric space and $\alpha>0$, then $(X,\alpha d)$ is a metric
space.  As $\alpha\to0$ large scales become visible, whereas as
$\alpha\to\infty$ small scales become visible.
\end{remark}

\begin{remark}[Quasi-metric spaces]
\label{rem:chow-iii-g2-quasi-metric}
\source{chow-et-al-ricci-flow-part-iii:page-0410}
\dcref{appG:G.2}\group{Metric spaces}
A map $\widetilde d:X\times X\to\R\cup\{\infty\}$ is a quasi-metric when
$\widetilde d(x,z)\leq\widetilde d(x,y)+\widetilde d(y,z)$.  A pseudo-metric
also satisfies nonnegativity, vanishing on the diagonal, and symmetry.  The
quotient by $x\sim y\iff\widetilde d(x,y)=0$ carries the induced metric
$d'([x],[y])=\widetilde d(x,y)$.
\end{remark}

\begin{proposition}[Pullback metric]
\label{prop:chow-iii-g3-pullback-metric}
\source{chow-et-al-ricci-flow-part-iii:page-0411}
\dcref{appG:G.3}\group{Metric spaces}
For a map $f:X\to(Y,d_Y)$, the pullback $f^*d_Y(x,x')=d_Y(f(x),f(x'))$
is a pseudo-metric on $X$; it is a metric when $f$ is injective.
\end{proposition}
\begin{proof}
\sketch
Nonnegativity, symmetry, and the triangle inequality are inherited from
$d_Y$.  If $f$ is injective, $f(x)=f(x')$ implies $x=x'$, so the separation
axiom also holds.
\end{proof}

\begin{definition}[Quasi-length space and associated distance]
\label{def:chow-iii-quasi-length-space}
\source{chow-et-al-ricci-flow-part-iii:page-0411,chow-et-al-ricci-flow-part-iii:page-0412}
\dcref{appG:G.4--G.5}\group{Length spaces}
A quadruple $(X,I,A,\cL)$ is a quasi-length space when $A$ is closed under
restrictions, concatenations, and linear reparametrizations, $\cL$ is additive
under concatenation, and partial lengths vary continuously.  Its associated
quasi-distance is
\[
d_{\cL}(x,y)=\inf\{\cL(\gamma):\gamma\in A,\ \gamma(a)=x,\ \gamma(b)=y\},
\]
with value $\infty$ when no admissible path joins $x$ to $y$.  If $\cL$ is
nonnegative, reparametrization invariant, and induces the topology of $X$,
then $(X,\cL)$ is a length space.  A rectifiable path admits an arc-length
parametrization $\beta=\widetilde\beta\circ\phi$ with $\phi$ continuous and
nondecreasing.
\end{definition}
\begin{proof}
\sketch
The concatenation axioms give the triangle inequality for $d_{\cL}$.
Reparametrization by the cumulative partial-length function gives the stated
arc-length factorization; constant subintervals are collapsed by $\phi$.
\end{proof}

\begin{proposition}[Metric induced by a length space]
\label{prop:chow-iii-g6-induced-metric}
\source{chow-et-al-ricci-flow-part-iii:page-0412}
\uses{def:chow-iii-quasi-length-space}
\dcref{appG:G.6}\group{Length spaces}

The associated distance $d_{\cL}$ of a quasi-length space is a quasi-metric.
For a length space it is a metric (possibly taking the value $\infty$); a
metric obtained this way is called intrinsic, and it is strictly intrinsic when
the length structure is complete.
\end{proposition}
\begin{proof}
\sketch
Concatenating paths through an intermediate point proves the triangle
inequality.  The length-space topology and nonnegativity give finiteness,
symmetry, and separation in the metric case.
\end{proof}

\begin{proposition}[Length induced by a metric]
\label{prop:chow-iii-g7-length-induced-by-metric}
\source{chow-et-al-ricci-flow-part-iii:page-0413}
\uses{prop:chow-iii-g6-induced-metric}
\dcref{appG:G.7}\group{Length spaces}

For a quasi-metric $(X,d)$, the variation
\[
\cL_d(\gamma)=\sup_{a=\tau_0\leq\cdots\leq\tau_k=b}
\sum_i d(\gamma(\tau_{i-1}),\gamma(\tau_i))
\]
is a quasi-length structure.  If $d$ is a metric, it is a length structure,
and $d_{\cL_d}=d$ whenever $d$ is intrinsic.
\end{proposition}
\begin{proof}
\sketch
Refining partitions proves additivity under concatenation and monotonicity of
partial variations.  The triangle inequality controls restrictions and linear
reparametrizations.  The intrinsic identity is the defining equality between
the metric and the infimum of lengths.
\end{proof}

\begin{example}[Riemannian manifolds]
\label{ex:chow-iii-g8-riemannian-length}
\source{chow-et-al-ricci-flow-part-iii:page-0413,chow-et-al-ricci-flow-part-iii:page-0414}
\dcref{appG:G.8}\group{Length spaces}
For a connected Riemannian manifold $(\mathcal M^n,g)$, piecewise $C^1$ paths
have length $\cL_g(\gamma)=\int|\dot\gamma|_g$, and
$d(x,y)=\inf_\gamma\cL_g(\gamma)$.  Completeness realizes this infimum by a
geodesic, so $(\mathcal M,d)$ is a metric length space.
\end{example}

\begin{definition}[Hausdorff and Gromov--Hausdorff distance]
\label{def:chow-iii-gromov-hausdorff-distance}
\source{chow-et-al-ricci-flow-part-iii:page-0414,chow-et-al-ricci-flow-part-iii:page-0415}
\uses{prop:chow-iii-g7-length-induced-by-metric}
\dcref{appG:G.7-gromov-hausdorff}\group{Gromov--Hausdorff convergence}

For subsets $A,B$ of $(Z,d_Z)$, define
$d_H^Z(A,B)=\inf\{\epsilon:A\subset N_\epsilon(B),B\subset N_\epsilon(A)\}$.
For compact metric spaces set
\[
\GH(X,Y)=\inf_{Z,f,g}d_H^Z(f(X),g(Y)),
\]
where $f,g$ range over isometric embeddings into $Z$.  This is a metric on
compact metric spaces.
\end{definition}

\begin{remark}[Gromov--Hausdorff metric properties]
\label{rem:chow-iii-g9-gh-properties}
\source{chow-et-al-ricci-flow-part-iii:page-0415}
\dcref{appG:G.9}\group{Gromov--Hausdorff convergence}
The Gromov--Hausdorff distance is nonnegative, symmetric, and satisfies the
triangle inequality; distance zero between compact spaces is equivalent to
isometry.
\end{remark}

\begin{definition}[Distortion and epsilon-isometry]
\label{def:chow-iii-g10-epsilon-isometry}
\source{chow-et-al-ricci-flow-part-iii:page-0415}
\uses{def:chow-iii-gromov-hausdorff-distance}
\dcref{appG:G.10}\group{Gromov--Hausdorff convergence}

The distortion of $f:(X,d_X)\to(Y,d_Y)$ is
$\operatorname{dis}f=\sup_{x,x'}|d_X(x,x')-d_Y(fx,fx')|$.
An $\epsilon$-net meets every point within $\epsilon$, and an
$\epsilon$-isometry has distortion at most $\epsilon$ and image an
$\epsilon$-net.
\end{definition}

\begin{remark}[Approximation maps need not be continuous]
\label{rem:chow-iii-g11-approximation-continuity}
\source{chow-et-al-ricci-flow-part-iii:page-0415}
\dcref{appG:G.11}\group{Gromov--Hausdorff convergence}
An epsilon-isometry is a set map; continuity is not part of its definition.
\end{remark}

\begin{lemma}[Gromov--Hausdorff distance and epsilon-isometries]
\label{lem:chow-iii-g12-epsilon-isometry}
\source{chow-et-al-ricci-flow-part-iii:page-0415,chow-et-al-ricci-flow-part-iii:page-0417}
\uses{def:chow-iii-g10-epsilon-isometry}
\dcref{appG:G.12}\group{Gromov--Hausdorff convergence}

If $\GH(X,Y)<\epsilon$, there are $2\epsilon$-isometries in both directions.
Conversely, an $\epsilon$-isometry implies $\GH(X,Y)<2\epsilon$.
\end{lemma}
\begin{proof}
\sketch
Embed the spaces into a common metric space realizing the Hausdorff distance
up to $\epsilon$.  Choosing points within $\epsilon$ gives maps with image
nets and distortion at most $2\epsilon$.  Conversely, the metric obtained by
gluing the two spaces along an $\epsilon$-approximation has Hausdorff distance
less than $2\epsilon$.
\end{proof}

\begin{remark}[Correspondences]
\label{rem:chow-iii-g13-correspondences}
\source{chow-et-al-ricci-flow-part-iii:page-0416}
\dcref{appG:G.13}\group{Gromov--Hausdorff convergence}
The Gromov--Hausdorff distance is exactly related to the distortion of
correspondences, up to the standard factor-two normalization.
\end{remark}

\begin{definition}[Pointed Gromov--Hausdorff distance]
\label{def:chow-iii-g14-pointed-gh}
\source{chow-et-al-ricci-flow-part-iii:page-0416}
\uses{def:chow-iii-gromov-hausdorff-distance}
\dcref{appG:G.14}\group{Gromov--Hausdorff convergence}

An $\epsilon$-approximate pointed isometry $f:(X,x_0)\to(Y,y_0)$ satisfies
\[
B_Y(y_0,\epsilon^{-1})\subset N_\epsilon(f(B_X(x_0,\epsilon^{-1}))),
\quad |d_Y(fx,fx')-d_X(x,x')|<\epsilon
\]
on $B_X(x_0,\epsilon^{-1})$, and sends $x_0$ to $y_0$.  The pointed
Gromov--Hausdorff distance is the infimum of such two-sided errors.
\end{definition}

\begin{lemma}[Pointed epsilon-isometries]
\label{lem:chow-iii-g15-pointed-epsilon}
\source{chow-et-al-ricci-flow-part-iii:page-0417}
\uses{def:chow-iii-g14-pointed-gh, lem:chow-iii-g12-epsilon-isometry}
\dcref{appG:G.15}\group{Gromov--Hausdorff convergence}

If $\GHpt((X,x_0),(Y,y_0))<\epsilon$, there are $2\epsilon$-isometries
$B_X(x_0,\epsilon^{-1})\to B_Y(y_0,\epsilon^{-1}+\epsilon)$ and conversely,
both preserving basepoints.
\end{lemma}

\begin{definition}[Precompactness]
\label{def:chow-iii-g16-precompactness}
\source{chow-et-al-ricci-flow-part-iii:page-0417}
\dcref{appG:G.16}\group{Gromov compactness}
A subset of a topological space is (sequentially) precompact if every sequence
has a subsequence converging in the ambient space.
\end{definition}

\begin{theorem}[Gromov precompactness, pointed Ricci bound]
\label{thm:chow-iii-g17-gromov-precompactness}
\source{chow-et-al-ricci-flow-part-iii:page-0417}
\dcref{appG:G.17}\group{Gromov compactness}
The collection of pointed complete $n$-manifolds with
$\operatorname{Ric}\ge(n-1)K$ is precompact for pointed
Gromov--Hausdorff convergence.
\end{theorem}

\begin{theorem}[Uniform total boundedness criterion]
\label{thm:chow-iii-g18-total-boundedness}
\source{chow-et-al-ricci-flow-part-iii:page-0417}
\dcref{appG:G.18}\group{Gromov compactness}
A family of compact metric spaces with uniformly bounded diameter and uniformly
bounded cardinalities of $\epsilon$-nets is precompact in Gromov--Hausdorff
distance.
\end{theorem}

\begin{theorem}[Pointed total boundedness criterion]
\label{thm:chow-iii-g19-pointed-total-boundedness}
\source{chow-et-al-ricci-flow-part-iii:page-0418}
\dcref{appG:G.19}\group{Gromov compactness}
If every pointed ball $B(x_0,\rho)$ in a family admits an $\epsilon$-net of
cardinality at most $N(\epsilon,\rho)$, then the family is pointed
Gromov--Hausdorff precompact.
\end{theorem}

\begin{definition}[Boundedly compact space and tangent cone]
\label{def:chow-iii-g20-tangent-cone}
\source{chow-et-al-ricci-flow-part-iii:page-0418}
\uses{def:chow-iii-g14-pointed-gh}
\dcref{appG:G.20}\group{Metric cones}

A metric space is boundedly compact if every closed bounded subset is compact.
When the limit exists independently of $\alpha_i\to\infty$, its tangent cone at
$p$ is
\[
(T_pX,d_p,0_p)=\lim_{\alpha_i\to\infty}(X,\alpha_i d,p).
\]
For a Riemannian manifold this cone is $T_p\mathcal M\cong\mathbf E^n$.
\end{definition}

\begin{proposition}[Rescaling tangent and asymptotic cones]
\label{prop:chow-iii-g23-g25-cone-rescaling}
\source{chow-et-al-ricci-flow-part-iii:page-0419}
\uses{def:chow-iii-g20-tangent-cone, def:chow-iii-g11-asymptotic-cone}
\dcref{appG:G.23,appG:G.25}\group{Metric cones}

If a tangent cone $(T_pX,d_p)$ or asymptotic cone $(AX,d_{AX})$ exists, then
$(T_pX,cd_p)$ or $(AX,cd_{AX})$ is isometric to the original cone for every
$c>0$.
\end{proposition}
\begin{proof}
\sketch
Replace the scaling sequence $\alpha_i$ by $c\alpha_i$ (or $\omega_i$ by
$\omega_i/c$); the two definitions then differ only by the displayed global
homothety.
\end{proof}

\begin{example}[Nonunique tangent and asymptotic limits]
\label{ex:chow-iii-g24-g28-nonunique-cones}
\source{chow-et-al-ricci-flow-part-iii:page-0419,chow-et-al-ricci-flow-part-iii:page-0420}
\dcref{appG:G.24,appG:G.28}\group{Metric cones}
Connected compact spaces can have two nonisometric tangent-type limits along
different blow-up sequences; connected noncompact planar subspaces can likewise
have two nonisometric asymptotic limits.  These exercises show that existence
and uniqueness in the definitions are genuine hypotheses.
\end{example}

\begin{definition}[Asymptotic cone]
\label{def:chow-iii-g11-asymptotic-cone}
\source{chow-et-al-ricci-flow-part-iii:page-0419}
\uses{def:chow-iii-g14-pointed-gh}
\dcref{appG:G.11-asymptotic-cone}\group{Metric cones}

When independent of the sequence $\omega_i\downarrow0$, the pointed limit
\[
(AX,d_{AX},0)=\lim_{\omega_i\to0}(X,\omega_i d,p)
\]
is the Gromov--Hausdorff asymptotic cone.  It is independent of the basepoint.
\end{definition}

\begin{example}[Elementary asymptotic cones]
\label{ex:chow-iii-g26-g27-asymptotic-examples}
\source{chow-et-al-ricci-flow-part-iii:page-0419,chow-et-al-ricci-flow-part-iii:page-0420}
\dcref{appG:G.26--G.27}\group{Metric cones}
Curvature-pinching cones and complete noncompact manifolds of zero aperture
have a half-line as asymptotic cone.  Hyperbolic space has no
Gromov--Hausdorff asymptotic cone.
\end{example}

\begin{definition}[Euclidean metric cone]
\label{def:chow-iii-g29-euclidean-cone}
\source{chow-et-al-ricci-flow-part-iii:page-0420}
\uses{def:chow-iii-g20-tangent-cone}
\dcref{appG:G.29}\group{Metric cones}

For $\diam(X,d)\leq\pi$, the Euclidean metric cone is the topological cone
with
\[
d_{\Cone(X)}([(x_1,r_1)],[(x_2,r_2)])=
\sqrt{r_1^2+r_2^2-2r_1r_2\cos d(x_1,x_2)}.
\]
Radial dilation by $c>0$ is a homothety.
\end{definition}
\begin{proposition}[The cone metric and Riemannian cone]
\label{prop:chow-iii-g30-g32-cone-metric}
\source{chow-et-al-ricci-flow-part-iii:page-0420,chow-et-al-ricci-flow-part-iii:page-0422}
\uses{def:chow-iii-g29-euclidean-cone}
\dcref{appG:G.30,appG:G.32}\group{Metric cones}

The formula in Definition~\ref{def:chow-iii-g29-euclidean-cone} is a metric.
For a Riemannian manifold $(M,g)$, the distance of $g_{\Cone}=r^2g+dr^2$ on
$M\times(0,\infty)$ agrees with the cone metric.
\end{proposition}
\begin{proof}
\sketch
The formula is the law-of-cosines distance in the Euclidean comparison sector;
the triangle inequality follows by concatenating sectors.  For the Riemannian
claim, minimizing the radial variational integral gives
$r(u)=(a\cos u+b\sin u)^{-1}$, whose length is precisely the law-of-cosines
expression.
\end{proof}

\begin{remark}[Cone motivation]
\label{rem:chow-iii-g31-cone-motivation}
\source{chow-et-al-ricci-flow-part-iii:page-0421}
\dcref{appG:G.31}\group{Metric cones}
The cone over the unit sphere is Euclidean space, and the metric
$g_{\Cone}=r^2g+dr^2$ produces the same formula by the Euler--Lagrange equation.
\end{remark}

\section{Aleksandrov spaces with curvature bounded below}

\begin{theorem}[Toponogov comparison]
\label{thm:chow-iii-g33-toponogov}
\source{chow-et-al-ricci-flow-part-iii:page-0423,chow-et-al-ricci-flow-part-iii:page-0424}
\dcref{appG:G.33}\group{Comparison geometry}
For a complete Riemannian manifold with sectional curvature at least $k$, every
geodesic triangle (within the positive-curvature perimeter restriction) has a
$k$-comparison triangle.  Its corresponding angles are no larger, the hinge
opposite side is no shorter, and points on corresponding sides satisfy
$d(q,s)\geq d(\widetilde q,\widetilde s)$.
\end{theorem}

\begin{remark}[Toponogov comparison variants]
\label{rem:chow-iii-g34-toponogov-variants}
\source{chow-et-al-ricci-flow-part-iii:page-0424}
\dcref{appG:G.34}\group{Comparison geometry}
Toponogov comparison also has hinge and monotonicity formulations; the latter
is recorded below as comparison-angle monotonicity.
\end{remark}

\begin{definition}[Aleksandrov space]
\label{def:chow-iii-g35-g36-aleksandrov}
\source{chow-et-al-ricci-flow-part-iii:page-0424,chow-et-al-ricci-flow-part-iii:page-0425}
\dcref{appG:G.35--G.36}\group{Comparison geometry}
A geodesic triangle in a complete length space consists of three points and
shortest paths joining them; a hinge is two such paths with a common vertex.
A complete locally compact length space has curvature $\geq k$ if every point
has a neighborhood on which comparison triangles satisfy the opposite-side
inequality.  The same inequality for every triangle defines curvature $\geq k$
in the large.
\end{definition}

\begin{theorem}[Toponogov globalization]
\label{thm:chow-iii-g37-globalization}
\source{chow-et-al-ricci-flow-part-iii:page-0425}
\uses{def:chow-iii-g35-g36-aleksandrov}
\dcref{appG:G.37}\group{Comparison geometry}

An Aleksandrov space of curvature $\geq k$ has curvature $\geq k$ in the large.
\end{theorem}

\begin{definition}[Comparison angle and direction]
\label{def:chow-iii-g38-g39-angle}
\source{chow-et-al-ricci-flow-part-iii:page-0425,chow-et-al-ricci-flow-part-iii:page-0426}
\uses{thm:chow-iii-g37-globalization}
\dcref{appG:G.38--G.39}\group{Comparison geometry}

For distinct $x,y,z$, the Euclidean comparison angle at $y$ is
\[
\widetilde\angle xyz=\cos^{-1}\frac{d(x,y)^2+d(y,z)^2-d(x,z)^2}{2d(x,y)d(y,z)}.
\]
In an Aleksandrov space these angles along two shortest paths are monotone as
the parameters increase, so their limit at the common initial point exists.
Thus every shortest path has a direction and the angle between two shortest
paths is well-defined.
\end{definition}
\begin{proof}
\sketch
The monotonicity is the hinge form of Toponogov comparison.  Taking the joint
limit in the monotone net gives the angle; taking the same path twice gives
zero.
\end{proof}

\begin{theorem}[Integer dimension and regularity]
\label{thm:chow-iii-g40-g42-regularity}
\source{chow-et-al-ricci-flow-part-iii:page-0426,chow-et-al-ricci-flow-part-iii:page-0427}
\dcref{appG:G.40,appG:G.42}\group{Aleksandrov structure}
The Hausdorff dimension of an Aleksandrov space is a finite nonnegative integer.
If its dimension is $n$, an open dense subset is homeomorphic to an
$n$-manifold.
\end{theorem}

\begin{definition}[Strainers]
\label{def:chow-iii-g41-strainers}
\source{chow-et-al-ricci-flow-part-iii:page-0427}
\uses{def:chow-iii-g35-g36-aleksandrov}
\dcref{appG:G.41}\group{Aleksandrov structure}

A point is $(m,\epsilon)$-strained if it admits pairs $(a_i,b_i)$ whose
comparison angles satisfy
$\widetilde\angle a_ipb_i>\pi-\epsilon$ and all cross-angles exceed
$\pi/2-10\epsilon$.  The pairs form an $(m,\epsilon)$-strainer.
\end{definition}

\begin{theorem}[Gromov--Hausdorff closure and compactness]
\label{thm:chow-iii-g43-g44-compactness}
\source{chow-et-al-ricci-flow-part-iii:page-0427}
\uses{def:chow-iii-g35-g36-aleksandrov, thm:chow-iii-g17-gromov-precompactness}
\dcref{appG:G.43--G.44}\group{Aleksandrov structure}

The Gromov--Hausdorff limit of Aleksandrov spaces of curvature $\geq k$ is a
complete length space of curvature $\geq k$.  The family of spaces with
dimension at most $n$ and diameter at most $D$ is Gromov--Hausdorff precompact.
\end{theorem}

\begin{theorem}[Splitting and diameter bounds]
\label{thm:chow-iii-g45-g46-splitting-diameter}
\source{chow-et-al-ricci-flow-part-iii:page-0428}
\uses{def:chow-iii-g35-g36-aleksandrov}
\dcref{appG:G.45--G.46}\group{Aleksandrov structure}

An Aleksandrov space of nonnegative curvature containing a line is isometric to
$Y\times\R$ for an Aleksandrov space $Y$ of nonnegative curvature.  If its
curvature is at least $k>0$, then $\diam X\leq\pi/\sqrt{k}$.
\end{theorem}

\begin{remark}[Volume comparison]
\label{rem:chow-iii-g47-volume-comparison}
\source{chow-et-al-ricci-flow-part-iii:page-0428}
\dcref{appG:G.47}\group{Aleksandrov structure}
Bishop--Gromov volume comparison extends to Aleksandrov spaces with curvature
bounded below.
\end{remark}

\begin{theorem}[Existence of the tangent cone]
\label{thm:chow-iii-g49-tangent-cone}
\source{chow-et-al-ricci-flow-part-iii:page-0429}
\uses{def:chow-iii-g29-euclidean-cone, prop:chow-iii-g30-g32-cone-metric, thm:chow-iii-g40-g42-regularity}
\dcref{appG:G.49}\group{Aleksandrov structure}

If $X$ is $n$-dimensional Aleksandrov with curvature $\geq k$, then $T_pX$
exists, has nonnegative curvature, and is isometric to the Euclidean cone over
the completed space of directions $\Sigma_p$.  For $n\geq3$, $\Sigma_p$ is
compact of curvature at least $1$; if $\Sigma_p\cong S^{n-1}(1)$, then $p$ has
a neighborhood homeomorphic to an open subset of $\R^n$.
\end{theorem}

\begin{definition}[Exponential map and cut locus]
\label{def:chow-iii-g50-exponential-cut-locus}
\source{chow-et-al-ricci-flow-part-iii:page-0429,chow-et-al-ricci-flow-part-iii:page-0430}
\uses{thm:chow-iii-g49-tangent-cone}
\dcref{appG:G.50}\group{Aleksandrov structure}

Let $\Omega_p$ consist of cone vectors $[(v,t)]$ realized by geodesics from
$p$.  Define $\exp_p([(v,t)])=x$ at the endpoint and let $C(p)$ be the cut
locus; $C_p\subset\Omega_p$ consists of vectors ending in $C(p)$.
\end{definition}

\begin{definition}[Stratification]
\label{def:chow-iii-g51-stratification}
\source{chow-et-al-ricci-flow-part-iii:page-0430}
\dcref{appG:G.51}\group{Aleksandrov structure}
A stratification $X=\bigsqcup_{i=1}^N X_i$ into topological manifolds has
$\dim X_1>\cdots>\dim X_N$ and closed unions
$X_k^+=\bigcup_{i=k}^N X_i$.  The sets $X_i$ are its strata.
\end{definition}

\begin{theorem}[Aleksandrov spaces admit stratifications]
\label{thm:chow-iii-g52-aleksandrov-stratification}
\source{chow-et-al-ricci-flow-part-iii:page-0430}
\uses{def:chow-iii-g51-stratification, thm:chow-iii-g49-tangent-cone}
\dcref{appG:G.52}\group{Aleksandrov structure}

Every finite-dimensional Aleksandrov space admits a stratification into
topological manifolds.
\end{theorem}

\begin{lemma}[First variation of distance]
\label{lem:chow-iii-g53-distance-derivative}
\source{chow-et-al-ricci-flow-part-iii:page-0431}
\uses{def:chow-iii-g35-g36-aleksandrov, def:chow-iii-g50-exponential-cut-locus}
\dcref{appG:G.53}\group{Semi-concavity}

If $\alpha$ is a minimal geodesic from $q$ to $p$, $\beta$ is unit speed from
$q$, and $p\notin C(q)$, then
\[
\left.\frac{d}{dt^+}\right|_{0}d(p,\beta(t))=-\cos\angle_q(\alpha,\beta).
\]
\end{lemma}

\begin{definition}[Semi-concavity and quasi-geodesics]
\label{def:chow-iii-g54-g56-semi-concavity}
\source{chow-et-al-ricci-flow-part-iii:page-0431,chow-et-al-ricci-flow-part-iii:page-0432}
\uses{lem:chow-iii-g53-distance-derivative}
\dcref{appG:G.54,appG:G.56}\group{Semi-concavity}

A locally Lipschitz $f$ is $\lambda$-concave when
$f\circ\gamma(s)-\lambda s^2/2$ is concave on every unit-speed geodesic; it is
semi-concave when this holds locally with some $\lambda$.  A path is a
quasi-geodesic if it preserves $\lambda$-concavity for every such function.
\end{definition}

\begin{example}[A non-semi-concave Lipschitz function]
\label{ex:chow-iii-g55-power-function}
\source{chow-et-al-ricci-flow-part-iii:page-0431}
\dcref{appG:G.55}\group{Semi-concavity}
For $\alpha\in[1,2)$, $f(x)=|x|^\alpha$ is Lipschitz near the origin but is
not semi-concave there.
\end{example}

\begin{theorem}[Perelman stability]
\label{thm:chow-iii-g57-stability}
\source{chow-et-al-ricci-flow-part-iii:page-0432}
\uses{thm:chow-iii-g43-g44-compactness}
\dcref{appG:G.57}\group{Aleksandrov structure}

For a compact $n$-dimensional Aleksandrov space $X$ of curvature $\geq k$, a
sufficiently close compact $n$-dimensional Aleksandrov space $Y$ of the same
lower curvature bound is homeomorphic to $X$.
\end{theorem}

\begin{remark}[Stability reference]
\label{rem:chow-iii-g58-stability-reference}
\source{chow-et-al-ricci-flow-part-iii:page-0432}
\dcref{appG:G.58}\group{Aleksandrov structure}
Kapovich gives a detailed expository account of Perelman's stability theorem.
\end{remark}

\begin{corollary}[Topology of small metric balls]
\label{cor:chow-iii-g59-small-balls}
\source{chow-et-al-ricci-flow-part-iii:page-0432}
\dcref{appG:G.59}\group{Aleksandrov structure}
For every point $p$ of a compact Aleksandrov space as in Theorem~\ref{thm:chow-iii-g57-stability},
some metric ball $B(p,\epsilon)$ is homeomorphic to the tangent cone $T_pX$.
\end{corollary}

\begin{theorem}[Regularity of Aleksandrov spaces]
\label{thm:chow-iii-g60-regularity}
\source{chow-et-al-ricci-flow-part-iii:page-0433}
\dcref{appG:G.60}\group{Aleksandrov structure}
There is a full-Hausdorff-measure subset $X_0$ whose complement contains the
singular set, and the metric on $X_0$ is induced by a $C^{1/2}$ Riemannian
metric extending continuously over $X-S_X$.  The space is approximately twice
differentiable almost everywhere.  A geodesically complete Aleksandrov space
is isometric to a $C^{1,\alpha}$ Riemannian metric on a smooth manifold.
\end{theorem}

\begin{remark}[Tangent-cone questions]
\label{rem:chow-iii-g21-g22-tangent-questions}
\dcref{appG:G.21--G.22}\group{Metric cones}\source{chow-et-al-ricci-flow-part-iii:page-0418,chow-et-al-ricci-flow-part-iii:page-0419}
For a boundedly compact subset $S\subset X$, it is natural to ask whether
$T_pS$ embeds isometrically in $T_pX$, and whether tangent cones exist for
locally convex subsets of Riemannian manifolds.
\end{remark}

\begin{remark}[Riemannian approximation problem]
\label{rem:chow-iii-g48-riemannian-limits}
\dcref{appG:G.48}\group{Aleksandrov structure}\source{chow-et-al-ricci-flow-part-iii:page-0428}
It is open whether every finite-dimensional Aleksandrov space of curvature
$\geq k$ is a limit of complete Riemannian manifolds with a uniform lower
sectional-curvature bound (possibly in higher dimensions).
\end{remark}

\begin{remark}[Further convergence questions]
\label{rem:chow-iii-g61-further-convergence}
\dcref{appG:G.61}\group{Aleksandrov structure}\source{chow-et-al-ricci-flow-part-iii:page-0433}
The behavior of these notions under Gromov--Hausdorff convergence and collapsing
is a further topic beyond this appendix.
\end{remark}
